\chapter{Closures}\label{chap:closures}

The model has more variables than equations. The closure chooses which
variables are fixed at their current values and which the model determines.
It expresses assumptions about how the economy adjusts to a shock, and results
can depend on it as much as on the shock itself. \manref{cge_closure} sets the
closure, and every simulation stores the closure under which it was solved.

\section{Closure options}\label{sec:closure-options}

A closure option fixes one variable of each pair and frees the other. The
default, which \manref{cge_calibrate} sets, is listed first for each block.

\begin{cgetable}{p{0.13\linewidth}p{0.29\linewidth}p{0.24\linewidth}p{0.20\linewidth}}
\toprule
Block & Option & Fixed & Free\\
\midrule\endhead
\bottomrule\endlastfoot
External & \cmd{flexexr}: flexible exchange rate & $\V{FSAV}$ & $\V{EXR}$\\
 & \cmd{fixexr}: fixed exchange rate & $\V{EXR}$ & $\V{FSAV}$\\[3pt]
Government & \cmd{flexgsav}: flexible government savings & $\V{TINSADJ}$, $\V{TQADJ}$, $\V{TMADJ}$, $\V{TVADJ}$, $\V{GADJ}$ & $\V{GSAV}$\\
 & \cmd{fixgsav} or \cmd{fixgsav direct}: direct tax rates adjust & $\V{GSAV}$, $\V{TQADJ}$, $\V{TMADJ}$, $\V{TVADJ}$, $\V{GADJ}$ & $\V{TINSADJ}$\\
 & \cmd{fixgsav sales}: sales tax rates adjust & $\V{GSAV}$, $\V{TINSADJ}$, $\V{TMADJ}$, $\V{TVADJ}$, $\V{GADJ}$ & $\V{TQADJ}$\\
 & \cmd{fixgsav import}: import tariffs adjust & $\V{GSAV}$, $\V{TINSADJ}$, $\V{TQADJ}$, $\V{TVADJ}$, $\V{GADJ}$ & $\V{TMADJ}$\\
 & \cmd{fixgsav vat}: VAT rates adjust & $\V{GSAV}$, $\V{TINSADJ}$, $\V{TQADJ}$, $\V{TMADJ}$, $\V{GADJ}$ & $\V{TVADJ}$\\[3pt]
Savings--investment & \cmd{savings}: savings-driven & $\V{MPSADJ}$ & $\V{IADJ}$\\
 & \cmd{investment}: investment-driven & $\V{IADJ}$ & $\V{MPSADJ}$\\[3pt]
Factor $f$ & \cmd{mobile}: fully employed and mobile & $\V{QFS}_f$, $\V{WFDIST}_{fa}$ & $\V{W}_f$, $\V{QF}_{fa}$\\
 & \cmd{specific}: fixed in each activity & $\V{W}_f$, $\V{QF}_{fa}$ & $\V{WFDIST}_{fa}$, $\V{QFS}_f$\\
 & \cmd{unemployed}: unemployment at a fixed wage & $\V{W}_f$, $\V{WFDIST}_{fa}$ & $\V{QFS}_f$, $\V{QF}_{fa}$\\
 & \cmd{wagecurve}, $f \in FL$ & $\V{WFDIST}_{fa}$, $\V{LSW}_f$ & $\V{W}_f$, $\V{QFS}_f$, $\V{QF}_{fa}$\\
 & \cmd{minwage}, $f \in FL$ & $\V{WFDIST}_{fa}$, $\V{LSM}_f$ & $\V{W}_f$, $\V{QFS}_f$, $\V{QF}_{fa}$\\[3pt]
Numeraire & \cmd{cpi} & $\V{CPI}$ & $\V{DPI}$\\
 & \cmd{dpi}: producer prices of domestic sales & $\V{DPI}$ & $\V{CPI}$\\
\end{cgetable}

The tax scaling factors $\V{TQADJ}$, $\V{TMADJ}$ and $\V{TVADJ}$ are fixed at
one, except the one that the government closure frees
(\secref{sec:replacement}). The rents $\V{RM}_c$, $\V{RE}_c$ and $\V{RA}_a$ are
fixed at zero under every closure; a constraint frees one of them, or another
variable, and adds an equation, so the system stays square
(chapter~\ref{chap:constraints}). Variables that do not exist for an element,
such as the imports of a commodity that is not imported, are neither fixed
nor free: they are left out of the system.

The closures differ in what adjusts. Under \cmd{flexexr} foreign savings are
fixed in foreign currency and the exchange rate moves the trade balance;
under \cmd{fixexr} the exchange rate is fixed and foreign savings absorb any
change in the current account. Under \cmd{flexgsav} tax rates and real
government consumption are fixed and government savings absorb changes in
revenue; under \cmd{fixgsav} government savings are fixed and one tax's rates
are scaled in proportion to keep them. Under \cmd{savings} the savings rates
are fixed and investment adjusts to available savings; under
\cmd{investment} real investment is fixed and the savings rates of all
domestic institutions are scaled to finance it. The closures of each factor
are described with its regime: fully employed factors that move between
activities at one average price, factors fixed in each activity that earn
activity-specific prices, and factors whose supply adjusts at a fixed price,
along a wage curve or above a minimum wage (chapter~\ref{chap:labour}).

\section{Swaps}\label{sec:swaps}

After the options above, a closure may be changed element by element, as
GAMS modellers do by fixing and freeing variables. A swap fixes one element of
a variable that the closure leaves free and frees one that it fixes, so the
number of free variables does not change. For example,
\begin{itemize}
\item \cmd{swap(GSAV GADJ)} fixes government savings and lets government
consumption adjust;
\item \cmd{swap(EXR CPI)} makes the exchange rate the numeraire;
\item \cmd{swap(QG(com_hlt) GADJ)} fixes government purchases of health
services, which can then be shocked, and scales the other purchases with
them.
\end{itemize}
Variables that the reduced system computes from their definitions, such as
household demand, cannot be fixed (\secref{sec:reduced}).

\section{Checks}

Every closure is checked when it is set. It must be square, with as many free
variables as equations. It must also leave every equation with at least one
free variable: a closure that leaves an equation without one makes the
Jacobian singular, and \manref{cge_closure} names the equation. These are
structural checks, over the pattern of the variables that each equation
contains.
